\section{Lexikální analýza}\label{sec:lexikalni-analyza}

Zdrojový soubor je pro počítač nejprve jen posloupností znaků. Syntaktický analyzátor by však pracoval velmi nepohodlně,~kdyby musel znovu a~znovu rozpoznávat,~které znaky tvoří jméno proměnné,~číslo nebo operátor. Tuto práci proto oddělíme.

\emph{Lexikální analyzátor} neboli \emph{lexer} čte znaky zdrojového programu a~seskupuje je do posloupnosti \emph{tokenů}. Každý token určuje druh rozpoznané části programu a~může nést také její hodnotu.

\begin{figure}[h]
    \centering
    \begin{tikzpicture}[node distance=22mm]
    \node[task,fill=blue!10,minimum width=38mm] (source)
        {zdrojový program};
    \node[task,fill=orange!12,minimum width=34mm,right=of source] (lexer)
        {lexer};
    \node[task,fill=green!10,minimum width=43mm,right=of lexer] (tokens)
        {posloupnost tokenů};
    \draw[diredge] (source)--(lexer);
    \draw[diredge] (lexer)--(tokens);
    \node[below=4mm of source,font=\small] {\texttt{position = speed * 60;}};
    \node[below=4mm of tokens,align=center,font=\small]
        {\texttt{id = id * num ;}};
\end{tikzpicture}

    \caption{Vstup a~výstup lexikálního analyzátoru.}
    \label{fig:lexer-vstup-vystup}
\end{figure}

\subsection{Lexémy a~tokeny}

\emph{Lexém} je konkrétní úsek zdrojového textu. \emph{Token} je informace,~kterou o~něm lexer předá další fázi. Pro příkaz
\[
    \texttt{position = speed * 60;}
\]
můžeme získat například následující tokeny:

\begin{center}
\renewcommand{\arraystretch}{1.2}
\begin{tabular}{|c|c|c|}
    \hline
    \textbf{lexém} & \textbf{druh tokenu} & \textbf{předaná hodnota}\\
    \hline
    \texttt{position} & identifikátor & \texttt{position}\\
    \texttt{=} & přiřazení & ---\\
    \texttt{speed} & identifikátor & \texttt{speed}\\
    \texttt{*} & násobení & ---\\
    \texttt{60} & číslo & $60$\\
    \texttt{;} & středník & ---\\
    \hline
\end{tabular}
\end{center}

U~identifikátoru potřebuje parser obvykle znát jeho text,~u~číselného literálu zase číselnou hodnotu. U~operátoru \texttt{+} často stačí samotný druh tokenu. Lexer si navíc obvykle ukládá pozici lexému,~aby později mohl vypsat srozumitelnou chybovou zprávu.

\subsection{Od regulárního výrazu k~lexeru}

Typické lexémy mají pravidelnou stavbu:
\begin{itemize}
    \item identifikátor může být popsán zkratkou
        \texttt{[A-Za-z\_][A-Za-z0-9\_]*};
    \item celé nezáporné číslo zápisem \texttt{[0-9]+};
    \item mezery a~konce řádků tvoří jazyk,~který lexer rozpozná,~ale nepředá parseru;
    \item každý pevný operátor,~například \texttt{+} nebo \texttt{=},~tvoří jednoprvkový jazyk.
\end{itemize}

Zápisy v~hranatých závorkách a~symbol \texttt{+} jsou běžné zkratky rozšířených regulárních výrazů: \texttt{[0-9]} znamená volbu jedné číslice a~\texttt{R+} jedno nebo více opakování výrazu \texttt{R}. Pomocí základních operací zavedených v~sekci \smartref{\showrefs}{sec:regex}{o~regulárních výrazech} je lze rozepsat pouze pomocí sjednocení,~zřetězení a~Kleeneho hvězdy.

Pro každý druh tokenu můžeme sestrojit konečný automat. Jejich počáteční stavy spojíme novým počátečním stavem pomocí $\lambda$-přechodů,~čímž získáme společný NFA. Pomocí \smartref[věty~]{\showrefs}{thm:subset-construction}{podmnožinové konstrukce} jej převedeme na DFA,~který pak čte zdrojový text znak po znaku.

\begin{figure}[h]
    \centering
    \begin{tikzpicture}[
    desc/.style={task,fill=blue!8,minimum width=37mm,font=\small},
    phasebox/.style={task,fill=orange!12,minimum width=23mm,font=\small}
]
    \node[desc] (id) at (0,1.3) {identifikátor\\\texttt{[A-Za-z\_][A-Za-z0-9\_]*}};
    \node[desc] (num) at (4.7,1.3) {číslo\\\texttt{[0-9]+}};
    \node[desc] (op) at (9.4,1.3) {operátory\\\texttt{+ - * / =}};

    \node[phasebox,minimum width=39mm] (nfa) at (4.7,-0.4)
        {společný NFA};
    \node[phasebox,minimum width=39mm] (dfa) at (4.7,-2.2)
        {ekvivalentní DFA};
    \node[task,fill=green!10,minimum width=31mm] (lexer) at (9.3,-2.2)
        {lexer};

    \draw[diredge] (id.south)--(nfa.north);
    \draw[diredge] (num.south)--(nfa.north);
    \draw[diredge] (op.south)--(nfa.north);
    \draw[diredge] (nfa)--node[midway,right=2mm,font=\small]
        {podmnožinová konstrukce}(dfa);
    \draw[diredge] (dfa)--(lexer);
\end{tikzpicture}

    \caption{Využití regulárních výrazů a~automatů při konstrukci lexeru.}
    \label{fig:lexer-automaty}
\end{figure}

\subsection{Nejdelší lexém a~priorita}

Při čtení vstupu může automat projít přijímajícím stavem a~později přijmout ještě delší část. Lexer proto obvykle používá pravidlo \emph{nejdelšího lexému}: pokračuje,~dokud lze získat delší platný lexém,~a~vrátí poslední přijatou možnost. Ze vstupu \texttt{abc123} tak vznikne jeden identifikátor,~nikoliv šest jednopísmenných tokenů.

Pokud stejný lexém vyhovuje více pravidlům,~rozhodne jejich priorita. Slovo \texttt{let} například splňuje obecný předpis identifikátoru,~ale chceme jej rozpoznat jako klíčové slovo. Podobně musí lexer před jedním znakem \texttt{=} upřednostnit dvouznakový operátor \texttt{==},~pokud jej jazyk obsahuje.

\warningbox{Lexer nerozhoduje,~zda tokeny tvoří správný program. Posloupnost \texttt{let + ; 7} lze bez problému rozdělit na tokeny,~ale její syntaktickou chybu odhalí až parser.}

\subsection{Implementace tokenů}

V~našem jazyce použijeme dva druhy příkazů a~výrazy s~identifikátory,~čísly a~čtyřmi aritmetickými operátory. Potřebujeme také tokeny pro přiřazení,~středník a~závorky. Jejich druhy a~samotný token zachycuje program \ref{lst:tokens}.

\lstinputlisting[
    caption={Druhy tokenů a~datová třída tokenu.},
    label={lst:tokens}
]{04-kompilatory/code/mini_compiler/tokens.py}

\subsection{Implementace lexeru}

Lexer v~programu \ref{lst:lexer} prochází vstup zleva doprava. Čísla a~identifikátory čte pomocí samostatných metod,~jednoznakové symboly vyhledává ve slovníku. Bílé znaky ignoruje a~text od \texttt{//} do konce řádku považuje za komentář.

\lstinputlisting[
    caption={Jednoduchý lexikální analyzátor.},
    label={lst:lexer}
]{04-kompilatory/code/mini_compiler/lexer.py}

Pro začátek programu \texttt{let x = 2 + 3 * 4;} vrátí lexer druhy tokenů
\[
\begin{aligned}
    &\texttt{LET},\ \texttt{IDENTIFIER},\ \texttt{EQUAL},\
      \texttt{NUMBER},\ \texttt{PLUS},\\
    &\texttt{NUMBER},\ \texttt{STAR},\ \texttt{NUMBER},\
      \texttt{SEMICOLON},\ \texttt{EOF}.
\end{aligned}
\]
Zvláštní token \texttt{EOF} označuje konec vstupu a~parseru usnadní rozhodování,~zda už přečetl celý program.
